\section{Discussion}
\label{sec:discussion}
The persistent negative gamma regime in our 2024 sample deserves emphasis. Historical options literature assumed regime-switching dynamics \cite{ni2005stock,garleanu2009demand}, where markets alternated between dealer long gamma (dampening volatility) and dealer short gamma (amplifying volatility). The 0DTE options explosion has fundamentally altered this dynamic \cite{cboe2023}, creating an environment where dealers remain structurally short gamma. Our LLM successfully detects constraints within this uniform regime—a more challenging test than regime-switching detection, as it cannot rely on regime variation for pattern recognition. The combination of high detection rate (69.4\%), significant predictive power (Granger p = 0.042), and 91.2\% materialization accuracy demonstrates a genuine understanding of structural mechanisms rather than statistical pattern-matching.

\subsection{Interpretation of Results}

Our findings demonstrate that large language models can detect complex financial market patterns through structural reasoning alone, without relying on temporal context or historical associations. The 71.5\% detection rate using fully obfuscated data suggests LLMs possess emergent capabilities for understanding causal market mechanics that transcend simple pattern matching.

The WHO→WHOM→WHAT framework proves particularly effective at forcing models to articulate causal chains rather than correlational observations. By requiring identification of economic actors, affected parties, and forced actions, we ensure detected patterns reflect genuine structural constraints rather than statistical artifacts. The consistent 91.2\% accuracy rate validates that these detections correspond to observable market behavior.

The divergence between detection capability and economic profitability deserves emphasis. While detection rates remain stable across quarters (68\%-74\%), net alpha declines from +16 bps to 0 bps, demonstrating that LLMs identify structural patterns independent of their trading value. This separation validates our methodology—we test understanding of market mechanics, not the ability to generate profits.

\subsection{Limitations}
Several limitations warrant consideration:
\textbf{Single Asset Focus:} Our analysis focuses exclusively on SPY, the most liquid options market globally, for methodological reasons that strengthen rather than limit validation. SPY concentrates dealer hedging activity, exhibits tightest bid-ask spreads minimizing GEX measurement noise, and serves as the primary 0DTE vehicle (85\%+ of same-day expiration volume). As a broad market index, SPY isolates pure dealer gamma constraints without confounding company-specific factors (earnings, sector rotation) that complicate causal attribution in individual equities. Demonstrating LLM detection in the most informationally efficient market provides conservative validation that should generalize to less efficient underlyings. However, cross-asset validation remains valuable future work, particularly testing whether detection persists in equity options with positive gamma regimes absent from our 2024 SPY sample.
\textbf{End-of-Day Granularity:} Daily snapshots may miss intraday dynamics critical to 0DTE hedging patterns.
\textbf{Simplified Gamma Calculations:} Standard Black-Scholes Greeks without smile adjustments may introduce noise in absolute GEX values.
\textbf{Reasoning Depth and Interpretability:} Like all transformer-based models, LLMs function as black boxes—we cannot trace the precise internal representations or attention mechanisms that produce pattern detections. This opacity limits our ability to explain why specific model weights encode gamma hedging knowledge or how the model arrived at a particular detection. However, our validation framework addresses this limitation through empirical verification: (1) obfuscation testing eliminates memorization pathways, (2) the WHO→WHOM→WHAT framework forces explicit causal attribution, (3) 91.2\% materialization rate validates output accuracy, (4) Granger causality tests confirm statistical predictive relationships (p=0.042), and (5) detection stability despite declining profitability demonstrates mechanism-based reasoning rather than profit optimization. While we cannot fully explain internal reasoning processes, we rigorously validate that detected patterns reflect genuine structural constraints. The 28.5 percentage point gap between biased and unbiased prompts quantifies prompt sensitivity but does not undermine the 71.5\% unbiased detection rate that exceeds our mechanical threshold.
\textbf{Temporal Limitation:} Our 2024 sample may not generalize to crisis periods or structural bull/bear markets.

\subsection{Implications for Financial Markets}
The ability to detect dealer constraints without historical context suggests LLMs could identify novel patterns in unexplored markets through pure structural reasoning. Risk managers could use LLM-detected patterns as early warning signals for volatility amplification, while the obfuscation testing methodology validates a new approach for distinguishing genuine structural constraints from historical correlations.

As LLMs increasingly influence trading decisions, regulators must understand their pattern detection capabilities. Our framework provides tools for assessing whether AI systems genuinely understand market dynamics or merely memorize historical patterns.

\subsection{Theoretical Contributions}

We introduce \textit{obfuscation testing} as a rigorous methodology for validating structural pattern detection by removing temporal context while preserving mechanical relationships. Our results demonstrate that transformer architectures exhibit emergent understanding of complex financial constraints, detecting dealer hedging patterns from pure gamma exposure values. The WHO→WHOM→WHAT framework proves effective for forcing explicit causal attribution in financial applications. Together, these contributions distinguish genuine causal reasoning from pattern memorization in LLM financial analysis.

The divergence between pattern detection and economic profitability (Figure~\ref{fig:quarterly_stability}) suggests LLMs identify structural regularities without engaging in entrepreneurial discovery. This distinction matters for deploying AI in financial markets: pattern recognition aids risk management and surveillance but differs from adaptive decision-making that generates trading edge.

\subsection{Future Directions}
Future work should validate detection across multiple asset classes and market regimes to establish generalizability. Intraday analysis with high-frequency data could reveal microstructure patterns invisible at daily granularity, particularly for 0DTE options. Testing ensemble methods across multiple LLMs could distinguish model-specific artifacts from robust pattern detection through consensus mechanisms.